\newpage{\ }
\cleardoublepage

\chapter[Capítulo 1. Introducción general]{Introducción general}\label{cap:intro}

% otra intro más grande, enfocada en el cambio global y en la ecología aplicada al manejo
En el contexto de cambio global, con una intensificación del uso humano de los ecosistemas, resulta necesario entender cómo los distintos impulsores de cambio alterarán la dinámica de los ecosistemas, y con ello, la vida de las sociedades que los habitan \citep{IPCC2023}. En este sentido, la ecología aplicada puede hacer contribuciones relevantes para el manejo. En particular, la ecología de paisajes puede ayudar a entender y predecir el funcionamiento de los ecosistemas a una escala regional, tomando en cuenta características espaciales idiosincráticas de cada región \citep{gergel2017learning,rego2019applied}. Este énfasis en el aspecto espacial de los procesos es particularmente necesario en algunos procesos intrínsecamente contagiosos, como los incendios \citep{keane2004classification, mckenzie2011landscape}. En ecosistemas terrestres, el fuego es uno de los principales estructuradores de la vegetación, y con ello, de la biodiversidad en general \citep{bond_global_2005, bond2005fire}. Por la estrecha relación entre el fuego y las condiciones atmosféricas, los regímenes de incendios se están viendo alterados en la mayor parte del mundo \citep{flannigan_implications_2009, flannigan_fuel_2016, ellis_global_2022}. De esta manera, resulta necesario evaluar y proyectar la actividad del fuego a escala regional, para así poder informar el manejo de los ecosistemas tomando en cuenta particularidades locales \citep{mckenzie2011landscape, moritz2012climate, bowman_vegetation_2020}. 

El entendimiento actual de los factores que controlan el fuego nos permite proyectar, al menos de forma cualitativa, cuáles serían las tendencias de cambio en los regímenes de incendios esperadas bajo el cambio climático en diversos ecosistemas del planeta. Los incendios requieren la combinación de fuentes de ignición, vegetación con suficiente biomasa y continuidad, que conforma el combustible, y condiciones atmosféricas que permitan que la vegetación se seque lo suficiente como para arder \citep{krawchuk_global_2009, krawchuk_constraints_2011, bowman2014pyrogeography}. De esta manera, en ecosistemas de baja productividad, como pastizales o matorrales abiertos, el fuego normalmente está limitado por la continuidad del combustible, mientras que en ecosistemas más productivos el factor limitante es la humedad del combustible \citep{pausas2012fuel, pausas_global_2013}. Los ecosistemas más productivos pueden enfrentar una mayor incidencia de incendios si el clima se vuelve más cálido y seco, ya que el combustible estará disponible para arder por períodos de tiempo más largos \citep{kitzberger_projections_2022}. En cambio, los sistemas de baja productividad sólo experimentarían más fuego con un aumento en las precipitaciones en algún momento del año seguido de una estación seca \citep{pricope2012spatio, fischer2015influence, villagra_spatial_2024}. De esta manera, la cantidad de combustible aumentaría y posteriormente perdería humedad, permitiendo la ocurrencia de incendios. Este marco conceptual puede orientar esfuerzos para predecir cambios en el régimen de incendios asociados al cambio climático en grandes escalas, desde continental a global, pero no es suficiente para hacer predicciones o proyecciones a escala regional (i.e., de decenas a unos cientos de km).

Los controles generales del fuego mencionados arriba aplican también a una escala regional, subcontinental, pero a menor escala cobran mayor importancia factores idiosincráticos de cada región, como la topografía, los rasgos de las especies de plantas dominantes, y la disposición en el paisaje de las fuentes de igniciones y de las barreras para el fuego (como lagos, zonas urbanas o roquedales), entre otros \citep{rollins2004mapping, curt2018modelling}. A su vez, desde el punto de vista del manejo, las predicciones cualitativas ofrecen limitada utilidad; distintas regiones pueden presentar características y tendencias cualitativamente similares en cuanto a la incidencia del fuego, pero aún así, pueden diferir notablemente en aspectos cuantitativos: aunque en sistemas con regímenes de incendios similares se espere, por ejemplo, un aumento de la incidencia del fuego, no es lo mismo esperar un aumento del 10 \% en el área quemada anual que un 200 \%. De esta manera, para anticiparnos a cambios en el régimen de incendios necesitamos un entendimiento más fino de los factores que regulan el fuego a escala regional, que no sólo represente una descripción cualitativa, sino más bien, cuantitativa \citep{bondavalli2009quantitative}.

El entendimiento y la predicción de sistemas ecológicos a escalas espaciales y temporales grandes (decenas a cientos de km, y en el orden de unas pocas décadas a siglos) rara vez pueden lograrse únicamente mediante experimentos u observaciones de campo o laboratorio. Esto ha llevado a la ecología de paisajes a utilizar modelos de simulación espacialmente explícitos para representar la dinámica de los ecosistemas, lo que nos permite evaluar las implicancias de nuestras hipótesis, contrastar patrones simulados con datos, e incluso predecir o proyectar cambios en los sistemas en función de alteraciones en sus controles \citep{turner2015landscape,gergel2017learning,rego2019applied}. La ecología del fuego utiliza este enfoque para generar proyecciones de la incidencia del fuego, así como para evaluar posibles transformaciones en el paisaje \citep{keane2004classification,mckenzie2011landscape}. Pero para desarrollar y aplicar estos modelos se necesita también modelar la dinámica de incendios, de manera espacialmente explícita, con un enfoque cuantitativo, y tomando en cuenta características regionales. El modelado espacialmente explícito del fuego se encuentra en un estado avanzado, existiendo modelos como FlamMap \citep{finney2006overview}, Prometheus \citep{tymstra2010development}, o Phoenix \citep{tolhurst2008phoenix}, que pueden predecir el avance del frente de un incendio con gran precisión. Al estar formulados en base a principios físicos, con un enfoque predominantemente mecanístico, estos modelos pueden ser utilizados en diversos ecosistemas, pero esta flexibilidad tiene un precio alto: requieren de abundante información muy costosa de obtener. Por ejemplo, modelos como FlamMap requieren mapas detallados del tipo de combustibles de la región a evaluar, como mapas de altura basal del dosel, profundidad del mantillo, carga de combustibles finos por clase de tamaño, y datos atmosféricos a resolución subdiaria. Si obtener ese tipo de información resulta inviable, se necesita desarrollar modelos que permitan proyectar cambios en el régimen de incendios utilizando información menos costosa de obtener o que ya se encuentra disponible, como mapas de vegetación, índices de vegetación derivados de imágenes satelitales, entre otros \citep{morales_stochastic_2015}. De esta manera, modelos con menor capacidad de generalización, formulados en base a las relaciones empíricas entre las variables disponibles (e.g., área quemada en función del tipo de vegetación) pueden resultar más adecuados. Pero a su vez, al carecer de una formulación mecanística, estos modelos empíricos requieren de un conocimiento acabado sobre los factores que controlan el fuego a nivel regional, para imponer restricciones que no están presentes en su formulación.

El estudio de los factores que regulan el fuego a escala de paisaje puede encararse con distintos enfoques. Por un lado, se pueden realizar experimentos de quema en laboratorio o en campo para evaluar cómo las características de la vegetación local afectan el comportamiento del fuego o el potencial de ignición \citep{butler2004radiation,bianchi_ignition_2014,mcallister2016critical}. Con un enfoque observacional y de menor costo, también se pueden comparar las características de la vegetación más relevantes para el fuego entre comunidades de plantas, por ejemplo, evaluando la continuidad y humedad del combustible y sus condiciones microclimáticas asociadas \citep{paritsis_pine_2018, tiribelli_changes_2018, bianchi_comparison_2018}. A su vez, esta información obtenida generalmente a escala de sitio puede combinarse con el estudio de patrones de área quemada en grandes extensiones, lo que permite poner en un contexto más amplio la evidencia obtenida a escala fina, y contribuye a detectar cuáles son los factores más condicionantes del fuego a escala de paisaje \citep{mermoz_deteccion_2002, holz_ecological_2012, paritsis_habitat_2013}. Con toda esta información local se pueden desarrollar o adaptar modelos de simulación de fuego que permitan proyectar cambios en el sistema relacionados al cambio climático. 

% intro patagonia

El norte de la Patagonia Andina es una de las regiones para las que se espera un marcado aumento en la incidencia del fuego en el siglo XXI debido al cambio climático \citep{kitzberger_projections_2022}. Dado que los bosques templados que allí habitan presentan muy baja resiliencia al fuego, y que presentan un gran valor ecológico y cultural, resulta de gran interés evaluar las posibles trayectorias del paisaje en las próximas décadas \citep{kitzberger_firevegetation_2016}. Para esto se requiere avanzar en el entendimiento de los factores que regulan el fuego en esta región y en el desarrollo de modelos de simulación de incendios para así, en el futuro, integrarlos en modelos dinámicos de vegetación. 

% Resumir más esto y conectarlo con el párrafo de arriba 

Los bosques andino-patagónicos presentan una elevada productividad, con un clima estacional que permite la ocurrencia de incendios en verano, la estación seca, momento en que la abundante biomasa vegetal presenta un bajo contenido de humedad \citep{veblen_fire_2003, veblen_historical_2008, kitzberger_firevegetation_2016}. El sistema está dominado principalmente por dos tipos de vegetación: bosques, relativamente poco inflamables y dominados por especies no rebrotantes, como lenga (\textit{Nothofagus pumilio}), coihue (\textit{Nothofagus dombeyi}) o ciprés de la cordillera (\textit{Austrocedrus chilensis}), y matorrales, más inflamables, dominados por especies rebrotantes, como ñire (\textit{Nothofagus antarctica}), caña colihue (\textit{Chusquea culeou}) y otras especies leñosas \citep{veblen_fire_2003}. Las especies que dominan los bosques son heliófilas, por lo que deben recolonizar un sitio quemado pocos años después de la ocurrencia del fuego ($<$ 20 años, variando entre especies), antes de que las especies rebrotantes que estaban en el sotobosque sombreen el sitio. Además, en esos años, requieren abundante precipitación para regenerar y la presencia de árboles semilleros en cortas distancias (entre 30 y 100 m, según la especie), por su baja capacidad \citep{cuevas_tree_2000, kitzberger_effects_2005, tercero-bucardo_field_2007, landesmann_importance_2018, landesmann_increased_2021}. Si no se cumplen estas condiciones, las especies rebrotantes de menor porte que se encontraban en el sotobosque y otras especies leñosas dispersadas por semillas pasan a dominar el sitio. En zonas de altitud baja o intermedia los bosques devienen en matorrales, mientras que a mayor altitud, en pastizales \citep{kitzberger_firevegetation_2016}. La combinación de la baja resiliencia de los bosques al fuego con su menor inflamabilidad genera una retroalimentación positiva entre fuego y vegetación \citep{kitzberger_decreases_2012, tepley_influences_2018}. Esto significa que el fuego promueve la retracción de los bosques, y el reemplazo de los bosques por comunidades más inflamables, como los matorrales, aumenta la  inflamabilidad del paisaje \citep{kitzberger_firevegetation_2016, tiribelli_non-additive_2019}. Esto resulta de gran relevancia, ya que los sistemas con retroalimentación positiva son altamente sensibles a cambios en sus forzantes, por ejemplo, el clima \citep{scheffer_alternative_2009, tepley_influences_2018}. En particular, con las tendencias climáticas regionales, hacia un clima más favorable para el fuego, se espera que los bosques se retraigan considerablemente en el futuro \citep{lloret2005fire, perry_explaining_2012, tepley_influences_2018, kitzberger_projections_2022}. 

% bosques vs matorrales: inflamabilidad contrastante, acá adentrarse

Uno de los factores que genera una retroalimentación positiva entre fuego y vegetación es la contrastante inflamabilidad entre bosques y matorrales. Los bosques presentan menor continuidad vertical, densidad y proporción de combustibles finos que los matorrales \citep{paritsis_positive_2015, blackhall_effects_2017, tiribelli_changes_2018}. A su vez, el dosel alto que presentan genera sotobosques con un microclima más húmedo y con menores extremos de temperatura \citep{paritsis_positive_2015}, lo que probablemente promueva una mayor humedad del combustible en comparación con los matorrales, con menor estructura vertical sombreando la superficie. Estas diferencias entre bosques y matorrales en cuanto a indicadores de inflamabilidad reportadas por estudios de campo coinciden con los patrones de ocurrencia de incendios a escala de paisaje: los bosques se queman con menor frecuencia y en menores extensiones que los matorrales \citep{mermoz_landscape_2005, morales_stochastic_2015, tiribelli_non-additive_2019}. 

% lo que sabemos, metido en modelos, y sus debilidades

El conocimiento sobre la dinámica de fuegos y vegetación en la región se encuentra representado en un modelo de simulación de incendios \citep{morales_stochastic_2015} y un modelo dinámico de vegetación (ver apéndices de \citealt{tiribelli_non-additive_2019} y de \citealt{kitzberger_projections_2022}). Sin embargo, estos modelos se encuentran en un estado de desarrollo incipiente, requiriendo mejoras para poder generar predicciones razonables. Estos modelos conciben el tiempo en forma discreta, con resolución anual, y el espacio, como una grilla a 60 m de resolución. La simulación de fuego se realiza con un modelo de igniciones (no publicado, pero ver apéndice de \citealt{kitzberger_projections_2022}) y un modelo de propagación \citep{morales_stochastic_2015}. El modelo de igniciones se ajustó con datos de focos del Parque Nacional Nahuel Huapi. El modelo simula un número anual de igniciones que varía según la precipitación. Permite simular igniciones por rayos, sin patrón espacial, e igniciones humanas, con mayor probabilidad de ocurrencia cerca de caminos. El modelo de propagación considera efectos del viento y la pendiente, la orientación de la ladera, y el tipo de vegetación, clasificado en bosque y matorral \citep{morales_stochastic_2015}. Dado que este modelo se ajustó utilizando datos de sólo nueve incendios, los autores no pudieron cuantificar efectos de las condiciones atmosféricas. Estos modelos se integran en un modelo de vegetación que simula sucesión y transiciones tras el paso del fuego \citep{kitzberger_projections_2022}. Sin embargo, para realizar proyecciones confiables, se requiere mejorar los módulos de fuego en algunos aspectos clave. 

Al ajustarse el modelo de propagación se estimó una probabilidad de quema notablemente menor en bosques que en matorrales, pero no se controló por factores físicos, como altitud y precipitación, que, estando correlacionados con el tipo de vegetación, pueden sobrestimar dichas diferencias en probabilidad de quema \citep{morales_stochastic_2015}. Si bien estas comunidades presentan características estructurales que explican una inflamabilidad contrastante, al estimar modelos en base a datos de ocurrencia de incendios (como en \citealt{morales_stochastic_2015}), la diferencia en incidencia del fuego entre bosques y matorrales también puede ser explicada por factores físicos. Por ejemplo, si los bosques son más comunes en zonas de mayor altitud y precipitación, su menor incidencia de fuego puede no deberse únicamente a sus características estructurales, sino también a las características climáticas de las zonas en que son más frecuentes. Esto implica que si en ambientes que presentan baja propensión al fuego por su clima el paso del fuego reemplaza a los bosques por matorrales, estos matorrales no serán tan inflamables como un matorral promedio, que ocurre en zonas climáticamente más favorables para el fuego (como baja altitud y menor precipitación). Otro factor limitante del modelo de \citet{morales_stochastic_2015} es que, habiéndose estimado en base a un acotado número de incendios, los autores no pudieron cuantificar efectos de las condiciones atmosféricas sobre la propagación del fuego, un aspecto fundamental para proyectar la dinámica de incendios bajo cambio climático. Por su parte, si bien el modelo de ignición considera las condiciones atmosféricas mediante la precipitación, no representa efectos de la vegetación y la topografía. 

% Lo que hicimos

El objetivo general de esta tesis fue mejorar nuestro entendimiento sobre los factores que regulan la incidencia del fuego en el norte de la Patagonia Andina y avanzar el desarrollo de modelos de simulación de fuego. Nuestro trabajo tuvo dos aristas, con énfasis en entender y en predecir o proyectar, respectivamente. En el primer caso, buscamos separar los efectos de la vegetación y variables físicas sobre la incidencia del fuego, utilizando evidencia a escala de sitio y a escala de paisaje. A escala de sitio, comparamos la humedad del combustible y las condiciones microclimáticas entre bosques y matorrales, utilizando sitios pareados para controlar factores abióticos que también pueden afectar la humedad del combustible (Capítulo~\ref{cap:palitos}). A escala de paisaje, analizamos los patrones espaciales de ocurrencia de incendios en un área de $\approx 29000 \ \text{km}^2$ y en un período de 24 años, evaluando en qué medida las diferencias en probabilidad de quema entre tipos de vegetación se explican por las características ambientales en donde ocurren o por diferencias intrínsecas, como estructura o humedad del combustible (Capítulo~\ref{cap:patrones}). Posteriormente, con un enfoque orientado a la predicción o proyección, nos propusimos volcar el entendimiento de los factores que regulan el fuego en un conjunto de modelos para simular el régimen de incendios, incluyendo la ignición y la propagación (Capítulo~\ref{cap:modelos}). A diferencia de esfuerzos previos, estos modelos se formularon con mayor resolución temporal, permitiendo caracterizar de forma más detallada el efecto de las condiciones atmosféricas sobre la incidencia del fuego. A su vez, modelamos los efectos de más tipos de vegetación (tres tipos de bosques, matorrales y pastizales), de variables físicas (topografía, precipitación), y de indicadores de influencia humana, con el desafío de utilizar la menor cantidad de parámetros posibles para mantener los modelos estimables con información limitada. Finalmente, utilizamos los modelos de simulación para proyectar el régimen de incendios en el futuro, bajo distintos escenarios de cambio climático.